\cotefile{1}
\scanpage{51}
% manuscript line 1
On en conclut le

% manuscript line 2
\uline{Théorème 1} (10) $E$ étant un espace de Hilbert, \deleted{\unsure{$\mu$ étant une}} \inserted{$\mu$ une mesure} positive sur un certain espace loc$^{\mathrm{t}}$ compact,
on a un isomorphisme canonique d'espaces hilbertiens

% manuscript line 5
\[
L^2(\mu)\overset{(2)}{\otimes}E=L^2_E(\mu) \tag*{(11)}
\]

% manuscript line 6
\uline{Corollaire 1} $I$ étant un ensemble quelconque d'indices, $E$
un \deleted{\unsure{certain}} espace de Hilbert, on a un isomorphisme
canonique

% manuscript line 8
\[
l^2(I)\overset{(2)}{\otimes}E=l^2_E(I) \tag*{(12)}
\]

% manuscript line 9
où le deuxième \unsure{mb|membre} est l'espace \deleted{\unsure{de toutes}} \inserted{des}
\deleted{\unsure{parties}} familles \inserted{de carré sommable} $(x_i)_{i\in I}$ d'éléments de $E$, \deleted{de carré} \inserted{et muni}
\deleted{\unsure{sommable}} \inserted{du} produit scalaire $((x_i),(y_i))=\sum(x_i,y_i)$.

% manuscript line 12
\uline{Corollaire 2} Soient $M,N$ deux espaces loc$^{\mathrm{t}}$ compacts,
munis de mesures positives $\mu,\nu$, Alors \deleted{\unsure{$E$}} on
a un isomorphisme canonique d'espaces hilbertiens

% manuscript line 15
\[
L^2(\mu)\overset{(2)}{\otimes}L^2(\nu)=L^2(\mu\otimes\nu) \tag*{(13)}
\]

% manuscript line 16
(où $\mu\otimes\nu$ est la mesure produit \inserted{de $\mu$ et $\nu$} sur $M\times N$)

% manuscript line 17
En effet, on a un isomorphisme canonique $L^2(\mu\otimes\nu)=L^2_{L^2(\nu)}(\mu)$

% manuscript line 18
En spécialisant $L^2(\mu)$ et $L^2(\nu)$ à $l^2(I)$ et $l^2(J)$, on
\deleted{\unsure{on}} retrouve la prop. 1.

% manuscript line 20
Soient maintenant $E$ et $F$ deux espaces de Hilbert,
en vertu de $E'=\bar E$, \deleted{\unsure{l'espace des}} il y a un
isomorphisme naturel de $\bar E\otimes F$ sur l'espace des
opérateurs \deleted{\unsure{mais}} continus de rang fini de
$E$ dans $F$, donné par

% manuscript line 25
\[
(\bar a\otimes b).x=(x,a).b\qquad! \tag*{(14)}
\]

% manuscript line 26
% archive sheet number: 52
\scanpage{52}
% manuscript line 1
\deleted{Montrons que \illegible{1} $\unsure{a\in E}$} Forme qui peut s'écrire, en \deleted{\unsure{considérant}} \deleted{\unsure{le vecteur}} \inserted{posant} $u=\bar a\otimes b$

% manuscript line 2
\deleted{\unsure{Maintenant on a de plus}}

% manuscript line 3
\[
\deleted{\unsure{(15)\quad\text{On a de plus}}}\qquad (u.x,y)=(u,\bar x\otimes y) \tag*{(15)}
\]

% manuscript line 4
\deleted{\unsure{On a|Comme}\ \illegible{1}} \deleted{\unsure{où dans le second}}
(dans le second \unsure{mb|membre}, on prend \deleted{\unsure{le}} produit scalaire
dans l'espace préhilbertien $\bar E\otimes F$), formule qui
\unsure{restera|reste} valable \inserted{par linéarité} pour \deleted{toute} $u\in\bar E\otimes F$. On en
conclut \deleted{que} $|(ux,y)|\leq\|u\|_2\,\|\bar x\otimes y\|_2=\|u\|_2\|\bar x\|\,\|y\|
=\|u\|_2\,\|x\|\,\|y\|$, d'où résulte que l'opérateur
défini par $u$ a une norme $\leq\|u\|_2$. Par
suite, l'application $u\unsure{\to|\mapsto}\tilde u$ de $\bar E\otimes F$ dans $L(E,F)$
se prolonge par continuité en une application linéaire
continue de \deleted{$\unsure{\bar E\otimes}$} norme $\leq 1$ de $\bar E\overset{(2)}{\otimes}F$. Cette
application est d'ailleurs biunivoque, car si $(e_i)_{i\in I}$ est
une base orthonormale dans $E$ \inserted{(donc aussi dans $\bar E$)}, $(f_j)_{j\in J}$ une base
orthonormale dans $F$, alors $\bar E\overset{(2)}{\otimes}F$ s'identifie
à l'espace \deleted{\unsure{des} \illegible{1}} des matrices infinies
$u=(u_{ij})_{ij\in I\times J}$ qui sont de carré sommable, \inserted{(prop. 1)} et
il est \deleted{\unsure{évident}} \inserted{trivial} que l'opérateur \deleted{\unsure{$E$}} défini par $u$
a \deleted{la} \inserted{pour} matrice $(u_{ij})$, donc $u$ est nul si cet
opérateur est nul. \deleted{Nous ne distinguerons pas par la
suite entre les $u\in\bar E\overset{(2)}{\otimes}F$ et les op. qu'ils définissent}

% manuscript line 23
\uline{Définition 2} Soient $E$ et $F$ deux espaces de Hilbert.
On appelle op. de H.-S. de $E$ dans $F$, \deleted{\unsure{toutes}} \inserted{les} applications
linéaires de $E$ dans $F$ définies par des éléments
de $\bar E\overset{(2)}{\otimes} F$. L'esp. de ces op., \deleted{\unsure{est}} muni de la structure

% manuscript line 27
% archive sheet number: 53
\scanpage{53}
% manuscript line 1
hilbertienne déduite de celle de $\bar E\overset{(2)}{\otimes}F$ par transport
de structure, est aussi noté $L^{(2)}(E,F)$, \inserted{ou $L^2(E)$ \deleted{\unsure{si}} si $F=E$.} La norme de \deleted{\unsure{l'op.}}
\deleted{\unsure{la} \illegible{1}} l'élément $u\in L^2(E,F)$ est notée $\|u\|_2$.

% manuscript line 4
\[
\text{On a }\deleted{\text{\unsure{Rappelons donc}}}\quad\boxed{\|u\|_\infty\leq\|u\|_2} \tag*{(16)}
\]

% manuscript line 5
Les op. de Hilbert-Schmidt de $E$ dans $F$
sont des applications compactes (car limites dans $L(E,F)$ d'él. de rang fini. Ce sont
\deleted{\unsure{tout de suite}} \deleted{ceux dont} \inserted{aussi exactement les op. dont} les matrices $(u_{ij})$ par rapport à des bases
orthonormales $(e_i),(f_j)$, \deleted{sont} de carré sommable
(matrices dites aussi de H.-S.) \inserted{et on a \inserted{alors} $\|u\|_2=(\sum|u_{ij}|^2)^{1/2}$.} Par la suite, nous
identifions $\bar E\overset{(2)}{\otimes} F$ et $L^2(E,F)$.

% manuscript line 7a
% slanted
\margin{\unsure{arranger} \deleted{\unsure{mettre le (\# 2)}} \uline{\unsure{Une proposition}}}

% manuscript line 11
\uline{\unsure{Théorème 2}} Soient $E$ et $F$ deux espaces de Hilbert,
soit $u$ une \deleted{op. de H.-S.} \inserted{application linéaire continue} de $E$ dans $F$
\deleted{et si $G$ est un \unsure{autre esp.} de Hilbert}

% manuscript line 14
1) Si $u$ est un op. de H.-S., il en est de $=$
de $Au$ (resp. $uA$) pour \deleted{\unsure{toute}} application linéaire
continue \deleted{de} \inserted{de $F$ dans} un espace de Hilbert $G$ \deleted{\unsure{dans}} (resp. de
$G$ dans $E$) \,; et on a

% manuscript line 18
\[
\deleted{(17)}\quad\boxed{\|Au\|_2\leq\|A\|\,\|u\|_2\quad
\text{(resp. }\|uA\|_2\leq\|A\|\,\|u\|_2\text{)}} \tag*{(17)}
\]

% manuscript line 19
2) Pour que $u$ soit un op. de Hilbert-Schmidt, il
faut et il suffit que $u^*$ soit un op. de
H.-S. (de $F$ dans $E$). \inserted{et alors $\boxed{\|u\|_2=\|u^*\|_2}$. Plus précisément,}
\deleted{On} a, pour $u,v\in L^{(2)}(E,F)$

% manuscript line 23
\[
\boxed{(u,v)=(v^*,u^*)} \tag*{(18)}
\]

% manuscript line 24
% slanted
\margin{\uline{Th. 3} (\unsure{Caractérisation spectrale})}

% manuscript line 24
3) Soit $h=\sqrt{u^*u}\,\deleted{\unsure{\in E}}$ \inserted{(op. hermitien positif dans $E$)}. Pour que $u$ soit un
H.-S., il faut et suffit \deleted{\unsure{que $h$ soit} \illegible{3}} \deleted{\unsure{que la suite} \illegible{2} $\unsure{\rho_i}$,}
\deleted{\unsure{des valeurs propres (i)} de $h$ soit de carré sommable,}
\deleted{\unsure{et alors}} que l'on ait $\sum\unsure{(\mathrm{v.s.})^2\,h|(\mathrm{v.s.})^2}<+\infty$. Alors

% manuscript line 28
\[
\boxed{\|u\|_2=\|h\|_2=(\sum\rho_i^2)^{1/2}} \tag*{(19)}
\]

% manuscript line 29
\deleted{\unsure{parmi les valeurs propres de} $\sqrt{u^*u}$}

% manuscript line 29a
% slanted
\margin{\uline{Corollaire 1} $u$ H.-S. $\Longleftrightarrow |u|$ H.-S. alors $\|u\|_2=\||u|\|_2$}

% manuscript line 30
\uline{Corollaire 2} $u$ H.-S. $\Longleftrightarrow u^*u$ Fredholm, et on a

% manuscript line 31
\[
\|u\|_2=\sqrt{\operatorname{Tr}\deleted{\unsure{(u)}}u^*u}=\sqrt{\|u^*u\|_1}
\]

% manuscript line 32
% archive sheet number: 54
\scanpage{54}
% manuscript line 1
\uline{Démonstration} \deleted{En effet} 1) Prouvons \unsure{p.ex.} que $Au$ est
un op. de Hilbert-Schmidt si $u$ \inserted{l'est, et que} \deleted{$\unsure{\|Au\|_2=}$}
$\|Au\|_2\leq\|A\|\,\|u\|_2$. En effet, soit $(e_i)$ une base orthonormale
dans $E$, on a $\deleted{\unsure{\sum\|ue_i\||\sum\|Aue_i\|}}\sum\underset{\unsure{?}}{\|Aue_i\|^2}\leq\deleted{\unsure{\|A\|}}$
$\sum\|A\|^2\|\deleted{\unsure{A}}ue_i\|^2=\|A\|^2\sum\|ue_i\|^2$, d'où \deleted{\unsure{il résulte que}} \inserted{notre assertion.}
L'assertion analogue pour $uA$ en résulte par transposition,
\deleted{\unsure{si} $\sum\|ue_i\|^2$ \unsure{est fini, il en est de même de}}
compte tenu de 2).

% manuscript line 9
\deleted{\unsure{$\sum$} \illegible{1}} \deleted{2) On a un \unsure{isom.} antilinéaire}
canonique de $\bar E\overset{(2)}{\otimes}F$ sur $\overline{\bar E\overset{(2)}{\otimes}F}$, espace
qui \deleted{peut} s'identifie à $\bar{\bar E}\overset{(2)}{\otimes}\bar F$ (formule 8)
$=E\overset{(2)}{\otimes}\bar F$, \deleted{\unsure{donc}} espace lui-$=$ \unsure{canoniquemt|canoniquement} isomorphe
à $\bar F\overset{\unsure{(2)|2}}{\otimes}E$. \deleted{\unsure{Il y a ainsi}} En d'autres termes, on a
un \deleted{\unsure{isom.}} isom. d'espaces de Hilbert de
$\deleted{\unsure{\overline{\bar E\otimes F}}}\ \overline{L^2(E,F)}$ sur $L^2(F,E)$, et on vérifie
aussitôt que \deleted{\unsure{cette}} cette application est donnée
par $\bar u\to u^*$ $(u\in L^2(E,F))$, d'où 2).

% manuscript line 18
3) On sait qu'on a $u=Uh$, où $U$ est
une application \unsure{partiellt|partiellement} isométrique de $E$ dans $F$,
et $h=Vu$, où $U$ est une application partiellement
isométrique de $F$ dans $E$. On en conclut, \deleted{que} \inserted{en}
\deleted{\unsure{si} \illegible{2} $u$} \inserted{vertu de 1),} que \deleted{\unsure{$h$ et $u$ sont}} \inserted{si l'un des} op. \deleted{\unsure{de}} $u$
et $h$ est de H.-S., il en est de $=$ de
l'autre, \deleted{et \unsure{de norme au plus égale}} \inserted{et a une norme de H.-S. au plus}
égale, d'où en \deleted{\unsure{effet}} \inserted{fait} $\|u\|_2=\|h\|_2$. \deleted{\unsure{On}} D'autre
\deleted{\unsure{est de}} \deleted{\unsure{si $h$ est}} \inserted{pour que $h$}
part, \deleted{pour que} soit un Hilbert-Schmidt, il
faut d'abord que $h$ soit cpt. Alors on sait

% manuscript line 29
% archive sheet number: 55
\scanpage{55}
% manuscript line 1
que, pour une base orthonormale convenable dans $E$,
$h$ se met sous forme d'une matrice diagonale \,;
\deleted{\unsure{(ses}} \inserted{alors} les \deleted{\unsure{termes}} diagonaux \deleted{sont} non nuls sont
\unsure{ses|les} valeurs propres \inserted{\unsure{non nulles} $(\rho_i)$} de $h$. On sait alors que $h$ est
de H.-S. si et seulement si cette matrice est
de carré sommable, i.e. si $(\sum\lvert\rho_i\rvert^2)^{1/2}<+\infty$, et
qu'on a alors $\|h\|_2=(\sum\rho_i^2)^{1/2}$. Le th. 2 est par là
complètement démontré.

% manuscript line 9
Les relations entre opérateurs de H.-S. et opérateurs
de Fredholm sont exprimées dans le

% manuscript line 11
\uline{Th. 3} Soient $E$ et $F$ deux espaces de Hilbert. L'application
linéaire canonique de $\bar E\widehat\otimes F$ dans $L(E,F)$ est biunivoque,
\deleted{\unsure{et}} Pour qu'une application linéaire $u$ de $E$ dans $F$
soit application de Fredholm, il faut et il suffit
que $h=\sqrt{u^*u}$ le soit. \deleted{\unsure{Pour}} Pour ceci, il faut et il suffit \unsure{que}
\deleted{et si les $(\lambda_i)$ sont les} $h$ soit compact et la \deleted{\unsure{suite}} $(\rho_i)$ de \unsure{ces|ses|les} valeurs propres
\deleted{et si les valeurs propres de $h$, \unsure{alors il est il suffit}
que $\sum\lambda_i<+\infty$, et alors on a} \deleted{$\unsure{\sum\lambda_i\geq0}$}
sommable, et alors on a

% manuscript line 20
\[
\|u\|_1=\|h\|_1=\sum\rho_i \tag*{(20)}
\]

% manuscript line 21
On peut aussi y remplacer $h$ par $\sqrt{uu^*}$.
\deleted{\unsure{Cela résulte}}

% manuscript line 22
% margin note written sideways
\margin{La validité des 2 premières \deleted{\unsure{parties}} assertions du théorème a déjà été signalée à la fin du \deleted{\unsure{au}} §2, N\textsuperscript{o} 1, (\uline{\unsure{parties en italiques|petites corrections}})}

% manuscript line 23
Que $u$ et $h$ soient simultanément des opérateurs
de Fredholm, et $\|u\|_1=\|h\|_1$, résulte du fait
qu'on a $u=Uh$, $h=Vu$, avec $\|U\|\leq1$, $\|V\|\leq1$.
\deleted{Si $h$ est \unsure{de}} \deleted{Si $h$ est tel \unsure{que} qu'il en est ainsi, et si on} Alors, $h$ sera un opérateur compact, donc
\deleted{$h=\sum\lambda_i\bar a_i\otimes a_i$} de la forme $\sum\rho_i\,\bar a_i\overset{\unsure{1|' }}{\otimes}a_i$, \deleted{où}
les $(a_i)$ forment un système orthonormal \inserted{dans $E$} et où les
$\rho_i$ sont les valeurs propres de $h$. Soit $p_n$ la
projection de $E$ sur l'espace engendré par les
$a_1,\ldots,a_n$, on aura $\sum_1^n\rho_i=\operatorname{Tr}p_nh\leq\|p_nh\|_1\leq\|h\|_1$
d'où $\sum\rho_i\leq\|h\|_1$. Réciproquement, supposons

% manuscript line 26a
% slanted
\margin{Non !}

% manuscript line 34
% archive sheet number: 56
\scanpage{56}
% manuscript line 1
$h$ compact et la série \inserted{$(\rho_i)$} de ses valeurs propres convergente, on aura
$h=\sum\rho_i\bar a_i\otimes a_i$, où $(a_i)$ est orthonormal \deleted{\unsure{est}}, d'où
$\|h\|_1\leq\sum\rho_i\|a_i\|^2=\sum\rho_i$, ce qui prouve que $h$ est op. de
Fredholm, et qu'on a (20). On peut d'ailleurs remplacer
\deleted{\unsure{le pr.}} $h=\sqrt{u^*u}$ par $h'=\sqrt{uu^*}$, \deleted{\unsure{cela se voit}} soit
directement, soit par application du fait \deleted{\unsure{suivant}} \inserted{bien connu} \deleted{que si $u$ est}
\deleted{\unsure{et $h$}} \uline{un op. compact quelconque} \inserted{$u\in L(E,F)$} \uline{alors} $\sqrt{u^*u}$ et $\sqrt{uu^*}$
\uline{ont les mêmes valeurs propres}. En effet, si $(\lambda_i)$ sont
les valeurs propres de $h=\sqrt{u^*u}$, on aura $h=\sum\rho_i\bar a_i\otimes a_i$
et $u=Uh$, où $U$ est isométrique sur l'espace engendré
par les $a_i$, d'où \unsure{$U$|$u$}$=\sum\rho_i\bar a_i\otimes b_i$, où $(b_i)=(Ua_i)$
forment un système orthonormal \inserted{dans $F$}. On en conclut donc

% manuscript line 13
\[
\begin{aligned}
u^*&=\sum\rho_i\bar b_i\otimes a_i,\quad\text{et}\quad uu^*=\sum_{ij}\rho_i\rho_j(a_i,a_j)\bar b_j\otimes b_i\\
&=\sum_i\rho_i^2\bar b_i\otimes b_i\,;\quad\text{d'où}\qquad
\sqrt{uu^*}=\sum_{\unsure{1|i}}\rho_i^2\bar b_i\otimes b_i,
\end{aligned}
\]

% manuscript line 15
\uline{Corollaire 1} Les op. de Fredholm de $E$ dans $F$ \deleted{\unsure{sont ceux de}} \inserted{sont \deleted{\unsure{ceux}} les}
opérateurs compacts de la forme

% manuscript line 16
\[
\boxed{u=\sum\rho_i\bar a_i\otimes b_i} \tag*{\textcircled{21}}
\]

% manuscript line 17
où $(a_i)$ (resp. $(b_i)$) est un système orthonormal
dans $E$ (resp. $F$), et $(\rho_i)\in l^1$. On a $\|u\|_1=\sum\rho_i$.

% manuscript line 19
% margin note written sideways
\margin{L'espace des applications de Fredholm de $E$ dans $F$ est noté $L^{(1)}(E,F)$, on a donc $L^{1}(E,F)=\bar E\widehat\otimes F=E'\widehat\otimes F$. \unsure{car $E'=\bar E$}.}

% manuscript line 19
Comme \unsure{$\bar E'$|$E'$}$\widehat\otimes E\to L(E)$ est biunivoque, on peut parler
de la trace d'un op. de Fredholm dans l'espace de Hilbert $E$.
\deleted{\unsure{On aura donc}} La formule (20) peut s'écrire aussi

% manuscript line 22
\[
\boxed{\|u\|_1=\operatorname{Tr}\sqrt{u^*u}=\operatorname{Tr}\sqrt{uu^*}} \tag*{\textcircled{22}}
\]

% manuscript line 23
\uline{Corollaire 2} Si $u$ est un op. de Fredholm de $E$ dans $F$,
on a

% manuscript line 25
\[
\|u\|_1=\sup_{\|A\|\leq1}|\operatorname{Tr} A\circ u|\qquad(A\in L(F,E)) \tag*{\textcircled{22}}
\]

% manuscript line 26
(En effet, le dual de \unsure{$E'$|$E''$}$\widehat\otimes F$ est $\unsure{B(E'',F)|B(E',F)}\simeq L(F,E)$) \inserted{\unsure{voir N. ant.}}

% manuscript line 27
% margin note written sideways
\margin{Conjuguant avec le th. 2, formule (19), on trouve le \uline{Corollaire 3} Pour que $u$ soit un op. de H.-S., il faut et il suffit que $u^*u$ soit un op. de Fredholm, et alors on a $(23)\quad\|u\|_2=\sqrt{\operatorname{Tr}(u^*u)}\deleted{\unsure{\leq}}=\sqrt{\|u^*u\|_1}$.}

% manuscript line 28
\deleted{\unsure{On}} Les relations entre opérateurs de Fredholm et op.
de H.-S. sont incluses dans l'important

% manuscript line 30
\uline{Th. 4} Soient \deleted{\unsure{$E$ et $F$, deux}} \inserted{$E,F,G$ des} espaces de Hilbert, $u\in L^2(E,F)$, $v\in L^2(F,G)$,
Alors $v\circ u\in L^1(E,G)$, et \deleted{\unsure{on}}

% manuscript line 32
\[
\boxed{\|vu\|_1\leq\|v\|_2\,\|u\|_2} \tag*{\textcircled{24}}
\]

% manuscript line 33
% archive sheet number: 57
\scanpage{57}
% manuscript line 1
De plus, si $u,v\in L^{(2)}(\deleted{\unsure{F}}E,\deleted{\unsure{E}}F)$, on a

% manuscript line 2
\[
\boxed{(u,v)=\operatorname{Tr}uv^*=\operatorname{Tr}v^*u} \tag*{\textcircled{25}}
\]

% manuscript line 3
\uline{Démonstration} Pour prouver (\unsure{25|24}), il suffit \deleted{\unsure{pour raison de continuité|par raison de continuité}}
\deleted{\unsure{si}} de se borner au cas où $u,v$ sont de rang fini
(alors \deleted{l'application} par prolongement par continuité on aura
une application bilinéaire de norme $\leq1$ de $L^2(E,F)\times
L^2(F,G)$ dans $\unsure{L^2(E,G)|L^1(E,G)}$, qui n'est \unsure{par} \deleted{\unsure{donc}} ne pourra
être que $(u,v)\to vu$). Or, \deleted{\unsure{ou|on}} d'après (22),
on est ramené \deleted{\unsure{al}} à prouver $|\operatorname{Tr} Avu|\leq\|v\|_2\|u\|_2$
pour \unsure{$\forall$} $A\in L(G,E)$, $\|A\|\leq1$. Comme $\|Av\|_2\leq\|A\|\|v\|_2$
\deleted{(prop. \unsure{1})} (formule (17)) on est ramené à prouver
$|\operatorname{Tr}wu|\leq\|w\|_2\|u\|_2$. En vertu de $\|v\|=\|v^*\|$, on
peut \deleted{\unsure{remplacer $v$ par $v^*$} (\illegible{1})} \inserted{poser $w=v^*$ \deleted{\unsure{$(v\in\bar E\otimes F)$}} $(v\in E'\otimes F)$}, et on est ramené
à prouver $|\operatorname{Tr}v^*u|\leq\|v^*\|_2\|u\|_2$. Mais cela résultera
de la formule (25), puisque $|(u,v)|\leq\|u\|_2\|v\|_2$.
\deleted{\unsure{Par|Ou}} Or la formule (25) pour opérateurs de rang
fini est triviale \,: il suffit de la vérifier pour
$u=\bar a_{\deleted{\unsure{1}}}\otimes b_{\deleted{\unsure{1}}}$, $v=\overline{a'}\otimes b'$, d'où $v^*=\overline{b'}\otimes a'$,
$v^*u=(b,\unsure{\overline{b'}|b'})\,\bar a\otimes a'$, \deleted{$\unsure{uv^*=(a,a')\bar b'\otimes b|uv^*=(a',a)\bar b\otimes b}$} d'où
\deleted{\unsure{Tr}} $(u,v)=(\bar a,\overline{a'})(b,b')=\deleted{\unsure{(a,a')|(a^*,a')}}\inserted{(a',a)}(b,b')$, $\operatorname{Tr}v^*u=(a',a)(b,b')$.
Comme $\operatorname{Tr}v^*u=\operatorname{Tr}uv^*$ (\inserted{formule de} commutation pour la trace)
(25) est démontrée pour $u,v$ de rang fini, \deleted{\unsure{Par là}} Ainsi
la première partie du th. est démontré, donc aussi
la formule (25) pour \deleted{tout} $u\in L^2(E,F)$, $v\in L^2(E,F)$, par
prolongement par continuité.

% manuscript line 26
% archive sheet number: 58
\scanpage{58}
% manuscript line 1
\uline{Corollaire 1} Pour que $u\in L(E,F)$ soit \unsure{op.} de H.-S., il faut et il
suffit que \deleted{$\unsure{uv}$} $vu$ soit op. de Fredholm pour \deleted{tout} $v\in L^2(F,E)$.
\deleted{Alors} \inserted{Alors}

% manuscript line 3
\[
\|u\|_2=\sup_{\|v\|_2\leq1}|\operatorname{Tr}vu|\qquad(v\in L^2(F,E)) \tag*{(26)}
\]

% manuscript line 4
En effet, on a vu que la condition est nécessaire, et
d'autre part la formule (25) s'écrit

% manuscript line 6
\[
\|u\|_2=\sup_{\|w\|_2\leq1}|\operatorname{Tr}w^*u|\qquad(w\in L^2(E,F))
\]

% manuscript line 7
ce qui résulte \unsure{immédiatement|aussi directement} de $\operatorname{Tr}w^*u=(u,w)$. Réciproquement,
si $vu\in L^1(E,E)$ pour \deleted{tout} $v\in L^2(F,E)$,
l'application $v\to vu$ de $L^2(F,E)$ dans $L^1(E)\deleted{\unsure{(E)}}$
est continue en vertu du th. du graphe fermé,
(car elle est continue quand $L^1(E)\deleted{\unsure{(E)}}$ est muni de
la topologie induite par $L(E)$, qui est séparée et moins fine).
Donc la forme linéaire $\bar w\to\operatorname{Tr}w^*u$ \deleted{\unsure{(sur)}} \inserted{sur} $\overline{\bar E\overset{(2)}{\otimes}F}$ est
continue, donc provient d'un élément $u_0$ de $\bar E\overset{(2)}{\otimes}F$
par la formule \deleted{\unsure{$\bar w\to(u,$}} $\operatorname{Tr}(w^*u)=(u_0,w)$ \,; faisant
$w=\bar x\otimes y$, on trouve que $(ux,y)=(u_0x,y)$ pour
\deleted{tout} $x\in E,y\in F$, d'où $u=u_0$\unsure{.}, d'où la conclusion.

% manuscript line 8a
% margin note written sideways
\margin{\unsure{supprimer ou mettre comme simple application|supprimer ou mettre comme remarque} }

% manuscript line 18
\uline{Corollaire 2} Pour que $u\in L(E)\deleted{\unsure{(E)}}$ soit op. de Fredholm,
il faut et il suffit que $u=vw$, où $v,w\in L^2(E)$, on
peut choisir $\|\unsure{u|w}\|_2=\|v\|_2=\sqrt{\|u\|_1}$.

% manuscript line 21
En effet, la suffisance \deleted{\unsure{est contenue}} \inserted{résulte} du th. 4, la
nécessité se voit en écrivant $u=Uh=Uw^2$, où
$w=\sqrt h$ \,; on a \unsure{manifestement} $\|w\|_2^2=\|h\|_1=\|u\|_1$,
et posant $v=Uw$, on aura $\|v\|_2\leq\|w\|_2$, et
\deleted{$\unsure{\|u\|_1}$} $u=vw$. On a d'ailleurs $\|v\|_2=\|w\|_2$ \inserted{$(=\sqrt{\unsure{\|u\||\|u\|_1}})$}, \deleted{\unsure{sinon}}
on a\deleted{\unsure{urait}} $\|u\|_1\leq\|v\|_2\|w\|_2<\|u\|_1$.

% manuscript line 27
% archive sheet number: 59
\scanpage{59}
% manuscript line 1
\uline{2. Produits extérieurs hilbertiens. Applications \unsure{à|:} des inégalités diverses}

% manuscript line 2
Soient $E_1,\ldots,E_n$ $n$ espaces de Hilbert. Alors on peut
encore, sur leur produit tensoriel $\bigotimes_{1\leq i\leq n}E_i$, introduire
un produit scalaire, donné sur les tenseurs décomposables
par

% manuscript line 6
\[
(a_1\otimes\cdots\otimes a_n,b_1\otimes\cdots\otimes b_n)=(a_1,b_1)\cdots(a_n,b_n) \tag*{(1)}
\]

% manuscript line 7
On vérifie comme au N° 1 qu'on obtient là une
forme sesquilinéaire \deleted{\unsure{pos}} strictement positive, \deleted{et que} d'où un
\deleted{$\unsure{\sum_{\alpha}\lambda_\alpha E_\alpha|\bigotimes_{i\in I_\alpha}^{(2)}}$ est une \unsure{base orthonormale, dont}}
\deleted{le complété de Hilbert \unsure{s'obtient} \illegible{1} noté}
\inserted{\uline{appelé \textquotedblleft produit tensoriel hilbertien des $E_i$\textquotedblright}, noté}
$\overset{(2)}{\bigotimes}_{1\leq i\leq n}E_i$, ou $\overset{n}{\otimes}\vphantom{\otimes}^{(2)}E$ si tous les $E_i$ se réduisent
à un $=$ espace $E$. Si, pour \deleted{\unsure{tout}} $i$,
$(e^i_{\alpha_i})_{\alpha_i\in A_i}$ est une base orthonormale dans $E_i$,
alors on obtient une base orthonormale
$(e_\alpha)_{\alpha\in A}$ de $\overset{(2)}{\bigotimes}_{1\leq i\leq n}E_i$, où $A=A_1\times\cdots\times A_n$,
et pour $\alpha=(\alpha_1,\ldots,\alpha_n)$, on pose $e_\alpha=e^1_{\alpha_1}\otimes\cdots e^n_{\alpha_n}$.

% manuscript line 9a
% margin note written sideways
\margin{\unsure{est un|on a un} isomorphisme canonique, \unsure{noté encore $\sigma$}, de $\overset{(2)}{\bigotimes}_{1\leq i\leq n}E_i$ sur $\overset{(2)}{\bigotimes}_{1\leq i\leq n}E_{\sigma\cdot i}$, \deleted{\unsure{isomorphisme} \illegible{2}}. Sur la formule (1), on voit aussitôt \uline{l'associativité} de l'opération du produit tensoriel hilbertien : si $1\leq\alpha<\beta\ldots<\gamma<n$, on a un isomorphisme canonique de $\overset{\unsure{(2)|n}}{\bigotimes}_{1\leq i\leq n}E_i$ avec $(\overset{(2)}{\bigotimes}_{1\leq i\leq\alpha}E_i)\overset{(2)}{\otimes}(\overset{(2)}{\bigotimes}_{\alpha<i\leq\beta}E_i)\overset{(2)}{\otimes}\cdots\overset{(2)}{\otimes}(\overset{(2)}{\bigotimes}_{\gamma<i\leq n}E_i)$. On voit de même qu'une \deleted{\unsure{autre}} permutation $\sigma\in\mathfrak S_n$, \unsure{sur} les tenseurs décomposables par $\sigma(x_1\otimes\cdots\otimes x_n)=x_{\sigma1}\otimes\cdots\otimes x_{\sigma n}$ \unsure{(commutativité du produit tensoriel hilbertien)}.}

% manuscript line 18
\deleted{Soit une} \inserted{Supposons donnée} une \deleted{\unsure{2ième}} deuxième \deleted{séquence}
$(F_i)_{1\leq i\leq n}$ d'espaces de Hilbert, et soit, pour \deleted{\unsure{tout}} $i$, \deleted{\unsure{$A$}} $u_i$ une application
lin. continue de $E_i$ dans $F_i$, alors l'application
$\bigotimes_{1\leq i\leq n}u_i$ de $\bigotimes E_i$ dans $\bigotimes F_i$ est continue
pour leurs normes hilbertiennes, et \uline{de façon précise}

% manuscript line 23
\[
\|u_1\otimes\cdots\otimes u_n\|\leq\|u_1\|\,\deleted{\unsure{\otimes}}\cdots\|u_n\| \tag*{(2)}
\]

% manuscript line 24
Il suffit en effet de le démontrer quand tous les $u_i$
sauf l'un sont l'identité. \deleted{On trouve le résultat en écrivant}

% manuscript line 26
% archive sheet number: 60
\scanpage{60}
% manuscript line 1
\deleted{\unsure{en groupant en produit, en écrivant} $u_1\otimes\cdots\otimes u$}
(car on a $u_1\otimes\cdots\otimes u_n=U_1\circ\cdots\circ U_n$, où $U_i$ est le
\deleted{$U_i=\unsure{1}\otimes\cdots u_i\otimes\cdots\otimes1$}
produit tensoriel de $n$ opérateurs, tous égaux à l'identité
sauf le $i$-ième, qui est égal à $u_i$. \deleted{Mais \unsure{identifiant les}}

% manuscript line 6
\deleted{Supposons donc $u\otimes1\otimes\cdots\otimes1$, où $A\in L(E_1,F_1)$,
identifiant $E_i$ à un espace $L^2(A_i)$, \unsure{avec} $F_i$ à un espace
$L^2(B_i)$, alors $\overset{(2)}{\bigotimes}_{1\leq i\leq n}E_i$ est identifié à l'espace $L^2(A_1\times\cdots\times A_n)$,
et $F_1\overset{(2)}{\otimes}E_{\unsure{2|1}}\overset{(2)}{\otimes}\cdots\overset{(2)}{\otimes}E_n$ est identifié à l'espace $L^2(A_1$}

% manuscript line 10
En vertu de l'associativité et de la commutativité du
produit tensoriel hilbertien, on \deleted{\unsure{peut}} est ramené alors au cas de
deux facteurs $E_1$ et $E_2$, $F_1$ et $F_2=E_2$, et où $A_1\in L(E_1,E_2)$, $A_2=\underline1$,
Mais alors, identifiant $E_1\overset{(2)}{\otimes}E_2$ à $L^{(2)}(\bar E_1',E_2)$, $F_1\overset{(2)}{\otimes}F_2$ à
$L^{(2)}(\bar F_1',E_2)$, \deleted{on est ramené} on est ramené
à prouver que pour $u\in L^2(\bar E_1',E_2)$, et $u$ de rang fini,
on a $\|u\,{}^tA_1\|_2\leq\|u\|_2\,\|\unsure{{}^tA_1|A_1}\|$, ce qui s'est en effet
avéré par la \deleted{formule de} deuxième formule (17) du N° 1,
aux notations près. \unsure{Par suite}, \deleted{l'indication \unsure{$u_1\otimes\cdots\otimes$}} \inserted{en vertu de (3),}
donc l'opération $\bigotimes_{1\leq i\leq n}u_i$ \deleted{\unsure{continue}} se \deleted{\unsure{transforme}} prolonge
par continuité en une opération de $\overset{(2)}{\bigotimes}_{1\leq i\leq n}E_i$ dans
$\overset{(2)}{\bigotimes}_{1\leq i\leq n}F_i$, notée encore de la $=$ façon. \deleted{\unsure{Reprenons les}}

% manuscript line 22
\deleted{\unsure{Soient donc donnés} des espaces de Hilbert $E_1,\ldots,E_n$.}
\inserted{Soit maintenant} $E$ un espace de Hilbert, considérons
$\overset{n}{\otimes}\vphantom{\otimes}^{(2)}E$, \deleted{Posons \unsure{encore}} Posons

% manuscript line 25
\[
a_n=\frac1{n!}\sum_{\sigma\in\mathfrak S_n}\varepsilon_\sigma\,\sigma \tag*{(3)}
\]

% manuscript line 26
où la somme du deuxième mb est considérée comme
somme d'opérateurs de symétrie dans $\overset{n}{\otimes}\vphantom{\otimes}^{(2)}E$ ($\varepsilon_\sigma$ signature
de la permutation $\sigma$). \deleted{\unsure{un proj}} On \unsure{reconnait|reconnaît}, au facteur $\frac1{n!}$

% manuscript line 29
% archive sheet number: 61
\scanpage{61}
% manuscript line 1
près, l'opérateur usuel d'antisymétrisation dans $\bigotimes_{1\leq i\leq n}E_i$. \deleted{\unsure{C'est}}
\deleted{\unsure{facile} \illegible{1} que $a_n$ est} $a_n$ est
évidemment un \uline{projecteur continu} (car il est continu, et
\deleted{\unsure{c'est un proj}} \inserted{on sait que} $a_n^2=a_n$ sur le \deleted{produit} produit tensoriel algébrique
en lui-même, qui est dense), d'ailleurs hermitien
(noter que l'adjoint de $\sigma$ est $\sigma^{-1}$, à cause de
la formule $(a_{\sigma1}\otimes\cdots\otimes a_{\sigma n},b_1\otimes\cdots\otimes b_n)=(a_1\otimes\cdots\otimes a_n,b_{\sigma^{-1}1}\otimes\cdots b_{\sigma^{-1}n})$
conséquence immédiate de (1)). \deleted{\unsure{La partie}}

% manuscript line 9
\uline{Définition} Soit $E$ un espace de Hilbert, $n$ un entier $\geq0$,
on appelle \deleted{produit} \inserted{puissance} \deleted{\unsure{ext.}} extérieure hilbertienne $n$-ième de $E$,
et on note $\overset{n}{\bigwedge}\vphantom{\bigwedge}^{(2)}E$, \deleted{\unsure{la}} \inserted{l'espace de Hilbert} quotient de $\overset{n}{\otimes}\vphantom{\otimes}^{(2)}E$ par
l'opérateur $a_n$ défini par (3), mais où le produit
scalaire est multiplié par $n!$.

% manuscript line 14
Cet espace s'identifie aussi à \deleted{un} un sous-espace
de $\overset{n}{\otimes}\vphantom{\otimes}^{(2)}E$ (image du projecteur $a_n$\inserted{(au facteur $n!$ près)}). Par définition,
on voit qu'il y a une application \uline{biunivoque}
naturelle du produit extérieur algébrique $\overset{n}{\bigwedge}E$ dans
$\overset{n}{\bigwedge}\vphantom{\bigwedge}^{(2)}E$, dont l'image est identique à l'image de $\overset{n}{\otimes}E$,
donc \uline{dense} dans $\overset{n}{\bigwedge}\vphantom{\bigwedge}^{(2)}E$. On identifie donc $\overset{n}{\bigwedge}E$ à \deleted{\unsure{un}}
son sous-espace préhilbertien de $\overset{n}{\bigwedge}\vphantom{\bigwedge}^{(2)}E$, donc ce dernier
est le complété de $\overset{n}{\bigwedge}E$. Le produit scalaire \deleted{dans}
$\overset{n}{\bigwedge}E$ s'explicite aisément \deleted{\unsure{après}} sur les \deleted{\unsure{mult.}} multivecteurs
décomposables : Posons $A=a_1\otimes\cdots\otimes a_n$, $B=b_1\otimes\cdots\otimes b_n$, on a

% manuscript line 24
\[
\deleted{(a_1\wedge\cdots\wedge a_n,b_1\wedge\cdots\wedge b_n)=\left(\frac1{n!}\right)^2\sum_{\substack{\sigma\in\mathfrak S_n\\\tau\in\mathfrak S_n}}\varepsilon_\sigma\unsure{\varepsilon_\tau a_1|\varepsilon_\sigma\,\sigma\,a_1}}
\]

% manuscript line 25
\[
\begin{aligned}
(a_1\wedge\cdots\wedge a_n,b_1\wedge\cdots\wedge b_n)&=n!(a_n A,a_n B)=n!(A,a_n^2B)=n!(A,a_nB)\\
&=\sum_{\sigma\in\mathfrak S_n}\varepsilon_\sigma(a_1,b_{\sigma1})\cdots(a_n,b_{\sigma n})=\det\big((a_i,b_j)\big),\quad\text{donc}
\end{aligned}
\]

% manuscript line 27
% archive sheet number: 62
\scanpage{62}
% manuscript line 1
\[
\boxed{(a_1\wedge\cdots\wedge a_n,b_1\wedge\cdots\wedge b_n)=\det\big((a_i,b_j)_{i,j\leq n}\big)} \tag*{(4)}
\]

% manuscript line 1a
% slanted, fragment on the folded corner
\margin{\unsure{(a)} \illegible{2} \unsure{(b)}}

% manuscript line 2
On peut donc énoncer la première partie du

% manuscript line 3
\deleted{\unsure{Prop.}} \uline{Théorème} Soit \deleted{\unsure{un}} $E$ un espace de Hilbert. Il existe sur
$(\overset{n}{\bigwedge}E)\times(\overset{n}{\bigwedge}E)$ une forme hermitienne définie positive (et une seule)
satisfaisant à 4. Le complété est la \deleted{produit tensoriel} \inserted{puissance extérieure} hilbertienne
$n$-ième de $\overset{n}{\bigwedge}E$. \deleted{\unsure{Si}} Si $(e_i)_{i\in I}$ est une base \deleted{\unsure{(hilbert}} orthonormée
de $E$, alors les $e_J$ donnés par (5) forment une base
\deleted{\unsure{hilbert. Notons que} \illegible{2} $(e_J)$ \illegible{3}} hilbertienne
\deleted{\unsure{Nous allons se} \illegible{2}} dans $\overset{n}{\bigwedge}\vphantom{\bigwedge}^{(2)}E$. Enfin, on a
\deleted{\unsure{on a} $\unsure{\det((a_i,a_j))}$}

% manuscript line 11
\[
\boxed{\|a_1\wedge\cdots\wedge a_n\|=(\det((a_i,a_j)))^{1/2}\leq\|a_1\|\cdots\|a_n\|} \tag*{(5)}
\]

% manuscript line 12
\deleted{En effet}, l'égalité \inserted{dans (5)} est un cas particulier de (4) (pour $b_i=a_i$).
\unsure{Considérons|Bornons-nous}, dans \deleted{l'espace engendré par les} \inserted{un espace de dim. $n$ contenant les} $a_i$, \deleted{une}
base orthonormale $(e_i)$, \deleted{telle que \unsure{l'on ait} \illegible{2}}
\deleted{vérifie aussitôt que} \inserted{posons} $a_{\unsure{i|j}}=\sum_j\deleted{\unsure{a_{ij}}}a_{ij}e_j$,
\deleted{\unsure{soit}} $A=(a_{ij})_{i,j\leq n}$, alors $(a_i,a_j)=\sum_k a_{ik}\overline{a_{jk}}$
\deleted{$\det((a_i,a_j))=\det$}
est l'élément d'indice $i,j$ de \inserted{la matrice} $AA^*$, donc
$\det((a_i,a_j))=\det AA^*=\det A\det A^*=\inserted{\det A\,\overline{\det A}=}|\det A|^2$, \deleted{\unsure{donc|d'où}} il
reste à prouver \fbox{(6) $\lvert\det A\rvert\leq\|a_1\|\cdots\|a_n\|$} (inégalité
de Hadamard \unsure{!}), qui exprime que le volume du
parallélépipède \unsure{rectangle|construit} sur les $a_i$ est au plus
égal au produit des longueurs des arêtes).

% manuscript line 15a
% slanted
\margin{\unsure{Ajouter au théorème pour préambule|Ajouter un théorème pour préambule} \illegible{2} que $X\to x\wedge X$ de $\overset{n}{\bigwedge}\vphantom{\bigwedge}^{(2)}E$ dans $\overset{n+1}{\bigwedge}\vphantom{\bigwedge}^{(2)}E$ est de norme $\leq\|x\|$.}

% manuscript line 24
Pour ceci, \inserted{il suffit de remarquer qu'}on peut supposer les $(e_i)$ choisis de façon
que $A$ soit une matrice triangulaire, donc
$\det A=a_{11}\,a_{22}\cdots a_{nn}$, ce qui achève la démonstration.

% manuscript line 27
\uline{Corollaire 1} Soient $a_i,b_i$ $(i=1,\ldots,n)$ des éléments d'un espace
de Hilbert $E$. Alors on a $|\det((a_i,b_j))|\leq\prod_i\|a_i\|\,\|b_i\|$

% manuscript line 29
% archive sheet number: 63
\scanpage{63}
% manuscript line 1
% Exposed upper edge of the underlying sheet. The slip covers the lower strokes.
\[
\boxed{(a_1\wedge\cdots\wedge a_n,b_1\wedge\cdots\wedge b_n)=\det\big((a_i,b_j)\unsure{,}\big.} \tag*{(4)}
\]

% manuscript line 2
% Beginning of the overlying slip.
Soit $(e_i)_{i\in I}$ une base hilbertienne \inserted{tot. ordonnée} de \deleted{\unsure{$F$}} $E$, pour \deleted{toute} partie
finie \inserted{de $n$ éléments} \deleted{$J\subset I$} $J$ de $I$, posons

% manuscript line 4
\[
\boxed{e_J=e_{i_1}\wedge\cdots\wedge e_{i_n}}. \tag*{(5)}
\]

% manuscript line 5
\inserted{$J=(i_1,\ldots i_n)$, $i_1<i_2<\cdots<i_n$,} On voit aussitôt sur (4) que
$(e_J,e_{J'})=\delta_{J,J'}$ (symbole de Kronecker), donc les
$e_J$ forment un système orthonormal dans $\overset{n}{\bigwedge}\vphantom{\bigwedge}^{(2)}E$. Ce système
est \deleted{évidemment} même \uline{total}, \inserted{car l'espace $E_0$ engendré par les $e_i$ est dense dans $E$, et} \deleted{car l'espace \unsure{engendré}} l'espace vectoriel engendré
\deleted{de $E_0^n$ dans $\overset{n}{\bigwedge}E$ par l'application multilinéaire}
par les $e_J$ \inserted{dans $\overset{n}{\bigwedge}\vphantom{\bigwedge}^{(2)}E$} est identique à celui engendré par l'image
de $(E_0)^n$ par l'application multilinéaire naturelle de $E^n$
\deleted{dans} sur $\overset{n}{\bigwedge}E$, donc elle est dense dans $\overset{n}{\bigwedge}E$ puisque
$E_0$ est dense dans $E$, donc dense dans $\overset{n}{\bigwedge}\vphantom{\bigwedge}^{(2)}E$.

% manuscript line 14
% archive sheet number: 64
% Exposed lower portion of the supporting sheet.

% manuscript line 15
\unsure{normale} $(e_i)$, \deleted{telle que \unsure{l'on ait} \illegible{2}}

% manuscript line 16
\deleted{vérifie aussitôt que} \inserted{posons} $a_{\unsure{i|j}}=\sum_j\deleted{\unsure{a_{ij}}}a_{ij}e_j$,
\deleted{\unsure{soit}} $A=(a_{ij})_{i,j\leq n}$, alors $(a_i,a_j)=\sum_k a_{ik}\overline{a_{jk}}$
\deleted{$\det((a_i,a_j))=\det$}
est l'élément d'indice $i,j$ de \inserted{la matrice} $AA^*$, donc
$\det((a_i,a_j))=\det AA^*=\det A\det A^*=\inserted{\det A\,\overline{\det A}=}|\det A|^2$, \deleted{\unsure{donc|d'où}} il
reste à prouver \fbox{(6) $\lvert\det A\rvert\leq\|a_1\|\cdots\|a_n\|$} (inégalité
de Hadamard \unsure{!}), qui exprime que le volume du
parallélépipède \unsure{rectangle|construit} sur les $a_i$ est au plus
égal au produit des longueurs des arêtes).

% manuscript line 18a
% slanted
\margin{\unsure{Ajouter} \unsure{préambule} \illegible{1} que $X\to x\wedge X$ de $\overset{n}{\bigwedge}\vphantom{\bigwedge}^{(2)}E$ dans $\overset{n+1}{\bigwedge}\vphantom{\bigwedge}^{(2)}E$ est de norme $\leq\|x\|$.}

% manuscript line 25
Pour ceci, \inserted{il suffit de remarquer qu'}on peut supposer les $(e_i)$ choisis de façon
que $A$ soit une matrice triangulaire, donc
$\det A=a_{11}\,a_{22}\cdots a_{nn}$, ce qui achève la démonstration.

% manuscript line 28
\uline{Corollaire 1} Soient $a_i,b_i$ $(i=1,\ldots,n)$ des éléments d'un espace
de Hilbert $E$. Alors on a $|\det((a_i,b_j))|\leq\prod_i\|a_i\|\,\|b_i\|$

% manuscript line 30
% archive sheet number: 63
\scanpage{64}
% manuscript line 1
En effet, il suffit, dans (4), de \unsure{s'écrire|s'aider de} l'inégalité de
Cauchy-Schwartz pour le premier mb, et d'utiliser (5).
Ce corollaire prouve que le coefficient $k_n$ qui intervient
dans Chap. 1, § 2, N° 2, prop. 1 est égal à 1 pour \deleted{\unsure{l'}} \inserted{un} espace de
Hilbert, (\deleted{\unsure{compte tenu}} \inserted{compte tenu} de l'identification \deleted{\unsure{du}} dual $E'$ de $E$ à $\bar E$)

% manuscript line 6
Soient maintenant $E,F$ deux espaces de Hilbert, et
soient $u_1,\ldots u_n$ $n$ applications linéaires continues
de $E$ dans $F$, considérons leur produit extérieur
$u_1\wedge\cdots\wedge u_n$ (Chap 1, § 1, N° 1). On a

% manuscript line 10
\[
u_1\wedge\cdots\wedge u_n=\deleted{\unsure{n!|1/n!}}\,\varphi_n\circ(u_1\otimes\cdots\otimes u_n)\circ a_n \tag*{(\unsure{8|7})}
\]

% manuscript line 11
où $a_n$ est l'isomorphisme \deleted{\unsure{naturel}} naturel de
$\overset{n}{\bigwedge}E$ dans $\overset{n}{\otimes}E$, et $\varphi_n$ l'homomorphisme naturel de
$\overset{n}{\otimes}F$ \deleted{\unsure{dans}} \unsure{sur} $\overset{n}{\bigwedge}F$. Comme $\frac1{\sqrt{n!}}a_n$ est \unsure{$=$|un}
isomorphisme de structures préhilbertiennes, et $\sqrt{n!}\,\varphi_n$ en est
de norme $\leq1$ pour les normes hilbertiennes mises
sur $\overset{n}{\otimes}F$ et $\overset{n}{\bigwedge}F$, on voit que

% manuscript line 17
\[
\|u_1\wedge\cdots\wedge u_n\|=\|\sqrt{n!}\,\varphi_n\circ(u_1\otimes\cdots\otimes u_n)\circ\frac1{\sqrt{n!}}a_n\|\leq\|u_1\otimes\cdots\otimes u_n\|,
\]

% manuscript line 18
d'où, en vertu de (2) :

% manuscript line 19
\[
\|u_1\wedge\cdots\wedge u_n\|\leq\|u_1\|\cdots\|u_n\|, \tag*{(9)}
\]

% manuscript line 20
Par suite, $u_1\wedge\cdots\wedge u_n$ peut être étendu en
une application linéaire continue de $\overset{n}{\bigwedge}\vphantom{\bigwedge}^{(2)}E$ dans
$\overset{n}{\bigwedge}\vphantom{\bigwedge}^{(2)}F$, \inserted{notée encore $u_1\wedge\cdots\wedge u_n$,} dont la norme satisfait à (9). Bien entendu,
les règles de calcul usuelles subsis\deleted{\unsure{tent}}tent pour les
\deleted{\unsure{produits}} $u_1\wedge\cdots\wedge u_n$. \inserted{(prolongés par continuité)} En particulier, si $u_i$ est un

% manuscript line 25
% archive sheet number: 65
